\section*{Comment 1 (Major) — Positioning and title}

\begin{quoting}
The manuscript is strongly positioned within mobile robotics, while the title and terminology are framed primarily around terrestrial laser scanning, which is commonly associated with surveying and geomatics. Please clarify whether the contribution is intended specifically for robotic platforms carrying survey-grade terrestrial laser scanners or for robotic stop-and-scan LiDAR mapping more generally. The terminology should then be applied consistently throughout the manuscript. A possible revised title can be “Occlusion-Aware Visibility Coverage for Robotic Stop-and-Scan 3D LiDAR Mapping”
\end{quoting}

\textbf{Response 1.} Thank you for pointing this out. We agree with this comment. Therefore, we have adopted the suggested title verbatim. The contribution is intended for robotic stop-and-scan LiDAR mapping in general; the survey-grade TLS (BLK360 G1) is the canonical instance of that sensor class and our experimental instrument. The revised Introduction now defines the regime explicitly (Section 1, p. 2): “We refer to this regime as \emph{robotic stop-and-scan LiDAR mapping}: a mobile platform carries a light detection and ranging (LiDAR) sensor that must remain stationary during each acquisition, and the dense scans are aligned offline rather than registered incrementally in real time. A terrestrial laser scanner (TLS) is the canonical instance of such a sensor … The method developed in this paper is formulated for the general stop-and-scan LiDAR regime; the BLK360 G1 enters only as the experimental instrument (Section 5).” Terminology has been made consistent throughout: the method sections (Sections 3–4) now speak of the generic stop-and-scan sensor with range bound R and stationary acquisition cost, and instrument-specific values appear only in the experimental sections (see also our response to Comment 5).

\section*{Comment 2 (Major) — Explicit research questions and a roadmap}

\begin{quoting}
The research gap is reasonably described, particularly in Section 2.4, where the authors identify the intersection between online exploration of unknown environments and planning for expensive stationary sensing. However, the manuscript does not formulate explicit research questions. The objectives and hypotheses must currently be inferred from the contributions and experimental sections. Please state one principal research question, supported where appropriate by subquestions. The Introduction should also conclude with a short roadmap explaining the structure of the remaining sections.
\end{quoting}

\textbf{Response 2.} We agree, and the Introduction now states one principal research question and three subquestions (Section 1, p. 3):
“\textbf{RQ.} \emph{In online robotic exploration of an unknown indoor environment with a stop-and-scan LiDAR sensor, does an occlusion-aware visibility coverage model produce better scan placement — higher genuinely observed coverage at a comparable number of stationary scans — than the conventional proximity-based disk model?}”
with subquestions (RQ1) how per-pose coverage should be defined on the live map so it never claims occluded or unmapped space (Section 3); (RQ2) how much of the placement difference is attributable to the coverage model alone, once exploration stochasticity \emph{and the acceptance rule} are controlled for (Sections 4 and 6); and (RQ3) whether the effect carries over to a physical platform and to the registered survey product (Sections 6.6 and 6.7). A structure roadmap paragraph now closes Section 1 (pp. 3–4).

\section*{Comment 3 (Major) — 2D visibility model vs. 3D point-cloud product}

\begin{quoting}
The proposed visibility model operates on a two-dimensional occupancy grid, whereas the resulting product and registration evaluation concern three-dimensional point clouds. The LOS metric therefore evaluates visible floor-space cells rather than the actual 3D surfaces observed by the scanner. It is unclear how the model accounts for vertical occlusions, shelves, elevated machinery, ceilings, overhanging structures, or surfaces that occupy the same horizontal position at different heights. High 2D LOS coverage does not necessarily imply comprehensive 3D surface coverage. The authors should either: extend the method toward a genuinely 3D visibility representation, such as a voxel-, octree-, or surface-based model; or explicitly define the contribution as 2D scan-station planning for subsequent 3D mapping, discuss this limitation, and identify full 3D visibility reasoning as future work. This distinction should be reflected consistently in the title, abstract, contributions, discussion, and conclusions.
\end{quoting}

\textbf{Response 3.} A fair and important point. We have taken the second option, and the 2D scope is now stated explicitly and consistently:
\begin{itemize}
\item \textbf{New Section 3.5 “Scope: 2D Station Planning for 3D Mapping” (p. 8)} states that the model reasons on a horizontal slice at sensor height; that LOS coverage certifies floor-space sight lines, not 3D surface coverage; that vertical occlusions (shelves, overhangs, elevated machinery, height-stacked surfaces) are not represented; why the 2D scope is matched to the online setting (the occupancy grid is what SLAM provides in real time; the 2D polygon is cheap enough to evaluate at every candidate; a full-dome scanner observes much of the vertical structure from a pose with clear 2D LOS); and that a genuinely 3D visibility representation (voxel-, octree-, or surface-based) is future work.
\item \textbf{Abstract (p. 1)}: “…a ray-cast visibility region computed on the robot’s live two-dimensional (2D) occupancy map … The planner thus performs 2D scan-station placement for subsequent 3D mapping.”
\item \textbf{Contributions (p. 3)}: the first contribution is now “defined on the live 2D occupancy grid”.
\item \textbf{Discussion (Section 7, p. 24)}: the first limitation is now the 2D representation and vertical occlusion, with 3D visibility identified as the principal future-work direction.
\item \textbf{Conclusions (Section 8, p. 24)}: “…the model performs 2D scan-station planning on the live occupancy map for subsequent 3D mapping”; future work includes the 3D extension.
\end{itemize}

\section*{Comment 5 (Major) — Method formulated independently of the instrument}

\begin{quoting}
The method should be formulated independently of the specific instrument used in the experiments. The Leica BLK360 G1 is an experimental implementation rather than an inherent component of the proposed algorithm. The BLK360 scan duration and accuracy should be presented as experimental settings rather than general characteristics of the proposed method.
\end{quoting}

\textbf{Response 5.} Agreed, and done. The method sections no longer reference the instrument: Section 4’s cost motivation now reads “minutes per acquisition for the survey-grade scanner used in our experiments (Section 5)” instead of quoting BLK360 scan times; the coverage-model figure caption (Figure 1, pp. 6–7) and the concept-figure title now refer to a generic scan pose; the summary and sweep figures label “stationary scans” rather than “BLK360 scans”. BLK360 duration, accuracy, and workflow details appear only in Section 5 (experimental setup, pp. 10–12) and in the hardware sections, as experimental settings.

\section*{Comment 6a (Major) — Broad claims, references, sensor grouping, abstract strength}

\begin{quoting}
Several broad claims require stronger references, particularly the statement that the coverage model using an isotropic disk is central to the scan-placement decision. Relevant literature on next-best-view planning, visibility analysis, view planning, and sensor placement should be cited. The manuscript groups structured-light scanners, thermal panorama systems, ground-penetrating radar, gas sensors, spectral sensors, and NDT probes as “high-precision stationary sensors.” These instruments vary in observation geometry, resolution, penetration, and uncertainty. They are better described as sensors needing stationary acquisition or dwell time. The current line-of-sight model is not directly applicable to penetrating sensors like GPR. In abstract, the statement that occlusion-aware visibility “is the appropriate basis” for stop-and-scan placement is too definitive given the limited environments, 2D representation, selected platform, and single scanner type.
\end{quoting}

\textbf{Response 6a.} We have made all three changes:
\begin{itemize}
\item \emph{Disk-model convention claim}: now supported with citations where first made (Section 1, p. 2) — “a sensor-footprint approximation widely used to drive online decisions in frontier exploration and coverage practice [Yamauchi 1997; Umari \& Mukhopadhyay 2017; Choset 2001; Galceran \& Carreras 2013; Placed et al. 2023].” The related-work section additionally cites the NBV/view-planning literature (Scott 2003; Tarabanis 1995; Zeng et al. 2020) and the visibility-based scan-planning literature (Mozaffar \& Varshosaz 2016; Prieto et al. 2017; Dehbi et al. 2021; Knechtel et al. 2025) (Sections 2.1–2.3, pp. 4–5).
\item \emph{Sensor grouping}: the family is now described as “sensors that require stationary acquisition or a dwell time at each viewpoint” (p. 2), and we state explicitly: “These instruments differ widely in observation geometry, resolution, and uncertainty; the line-of-sight coverage model developed here applies to the optical, non-penetrating members of this family and explicitly not to penetrating modalities such as GPR (Section 7).”
\item \emph{Abstract claim}: softened and scoped (p. 1) — “The results indicate that, for the room-scale indoor environments studied, occlusion-aware visibility is a sounder basis than Euclidean proximity for stop-and-scan placement.” The Conclusions carry the same scoping (“Within the room-scale indoor environments and the optical stop-and-scan sensor class studied here…”, p. 24).
\end{itemize}

\section*{Comment 6b (Major) — Conference-paper overlap and disclosure}

\begin{quoting}
The manuscript states that it extends a conference paper that is currently under review. Please verify that this is consistent with the journal’s editorial policy and clearly disclose the overlap between the two submissions. The authors should identify which methods, experiments, figures, and results are new in the journal submission. The sentence stating that the mobile platform was absent from the conference version is unnecessary in the scientific narrative and should be removed.
\end{quoting}

\textbf{Response 6b.} The Introduction now discloses the overlap explicitly and itemizes what is new (Section 1, p. 3):
“The conference version contributes the stop-and-scan system integration (frontier exploration coupled to a scan sequencer) and an isotropic-disk scan trigger, demonstrated in simulation. New to this article are: the ray-cast visibility coverage model and the LOS coverage metric (Section 3); the dual-criterion marginal-gain placement rule and the coverage-completion phase (Section 4); the controlled paired ablation with its three baselines and the parameter sensitivity analysis (Section 6); and the entire hardware study — the fully autonomous on-hardware disk-versus-visibility comparison and the per-run offline registration of the acquired scans (Sections 6.6 and 6.7). No figure, table, or quantitative result is shared between the two manuscripts.”
The sentence noting the platform’s absence from the conference version has been removed, as requested. Since the original submission, the conference paper has been \textbf{accepted for presentation at ICCAS 2026} (Regular Paper; decision of 31 July 2026); its status has been updated in reference [9], in the Introduction, and in the extended-version footnote on the first page. The accepted conference manuscript has been provided to the editorial office for the editors and reviewers, and we have confirmed that this extension route is consistent with MDPI’s editorial policy on extended conference papers.

\section*{Comment 7 (Major) — Equations (1)–(7) clarifications; disk baseline vs. Equations (6)–(7)}

\begin{quoting}
Equations (1)–(7) are generally understandable, but several definitions should be clarified: Define the grid-cell domain, cell-center coordinates, grid resolution, and ordering of the occupancy thresholds. Distinguish clearly between a grid cell and a candidate scan pose; the symbol (c) is probably used for both. Explain whether visibility regions associated with previous scan poses are frozen or recomputed when the online occupancy map changes. Describe Equation (6) as a normalized marginal gain and define it only for a nonempty visibility region. State the square-cell assumption required by Equation (7) and report the numerical grid resolution. The relationship between the disk baseline and Equations (6)–(7) also needs clarification. Skipping a candidate because its center lies within (R) of an existing scan is not mathematically equivalent to applying the marginal-gain rule B\_disk. If the baseline uses a different acceptance rule, the comparison changes both the coverage model and the decision rule.
\end{quoting}

\textbf{Response 7.} All requested clarifications have been made in Sections 3–4:
\begin{itemize}
\item \textbf{Grid domain/resolution/ordering} (Section 3.1, p. 6): “The grid is an H×W array of square cells of side ρ, the grid resolution (ρ = 0.05 m in all experiments); each cell c is identified with its center point in the map frame, so that ‖c − c′‖ denotes the Euclidean distance between cell centers. … thresholds ordered as 0 ≤ θ\_free < θ\_occ ≤ 100.”
\item \textbf{Cell vs. candidate pose} (Section 3.1, p. 6): the candidate pose is now denoted \textbf{p} throughout Section 4 and Algorithm 1 (“Throughout the paper, s\_i denotes an accepted scan pose and p a candidate pose under evaluation, while c is reserved for grid cells”).
\item \textbf{Frozen or recomputed} (Section 4.1, p. 8): “The visibility regions of previously accepted scans are not frozen: at every decision, C(S) is recomputed on the current occupancy map, so earlier regions expand or contract as SLAM refines the map (the final reported coverage is evaluated once more on a single reference map; Section 5).” This matches the implementation.
\item \textbf{Equation (6)} (p. 8): now introduced as the “\emph{normalized marginal gain}”, “defined only for a nonempty visibility region”, with the degenerate-region cut stated to reject candidates \emph{before} the ratio is evaluated, “so the denominator of the normalized gain never vanishes.”
\item \textbf{Equation (7)} (pp. 8–9): “where cells are squares of side ρ (the grid resolution, ρ = 0.05 m), so each cell contributes area ρ² and a is in m².”
\item \textbf{Disk baseline vs. Equations (6)–(7)}: the reviewer is right, and this observation prompted the main new experiment of the revision. We now state explicitly that the spacing rule is \emph{not} mathematically equivalent to the marginal-gain rule applied to B\_disk (“a candidate closer than R to an existing scan can still contribute a crescent of new disk area”, Section 4.2, p. 9), and we resolve the resulting confound with a \textbf{new disk marginal-gain baseline}: the identical dual criterion of Equation (8) with B\_disk substituted for B, replayed on the same fixed candidate paths (Sections 4.2 and 6.2; Table 2, p. 16; Figure 8, p. 19). Within that pair, the only difference is the coverage model. Results: single-room 77.7 ± 3.8\% LOS at 2.6 ± 0.5 scans (essentially the spacing rule’s result); multi-room \textbf{59.1 ± 1.3\% LOS} at 2.9 ± 0.3 scans — \emph{worse} than the spacing rule, because the disk model’s own coverage bookkeeping saturates through partitions and stops scanning while whole rooms remain unobserved (mechanism explained in Section 6.2, pp. 15–16). This isolates the deficiency to the coverage model itself rather than the acceptance rule.
\end{itemize}

\section*{Comment 8 (Major) — Figure 11 (parameter sweep) presentation}

\begin{quoting}
The general trends in Figure 11 are understandable, but the presentation and explanation need clarification. Please: identify the solid LOS-coverage curves and dashed scan-count curves clearly in the legend; avoid describing scan-count changes in “percentage points”; explain how (A\_min=5 m2) was identified as the knee point; and include variability measures for the averaged results.
\end{quoting}

\textbf{Response 8.} The figure has been regenerated (now Figure 8, p. 19, after the renumbering of Comment 13) and the text rewritten:
\begin{itemize}
\item The legend now identifies each curve explicitly: “single-room — LOS coverage (solid, left axis)”, “single-room — scan count (dashed, right axis)”, etc.
\item All curves now carry \textbf{error bars (per-path standard deviation over the N = 5 / N = 10 replay paths)}; the caption states “mean ± s.d.”
\item Scan-count changes are no longer described in “percentage points”: “…changes mean LOS coverage by at most about two points (single-room 86.4→84.3\%; multi-room 85.2→84.0\%) and mean scan count by less than one scan.”
\item The knee identification is now an explicit criterion: “it is the largest threshold whose mean LOS remains within about five points of the most aggressive setting (A\_min = 1 m²) in both environments (85.4 vs. 90.7\%; 84.0 vs. 85.7\%) while already halving the scan count (3.6 vs. 7.2 scans in both environments). Beyond the knee the trade turns unfavorable: pushing to A\_min = 15 m² saves only about one further scan (3.6→2.2 and 3.6→2.8) while coverage continues to erode (85.4→80.8\% and 84.0→83.1\%).”
\end{itemize}

\section*{Comment 9 (Minor) — Geometry vs. RGB}

\begin{quoting}
Please distinguish between the geometric measurements produced by the laser scanner and the RGB information supplied by the integrated camera. Unless color is used in visibility evaluation or registration, repeated references to “colorized point clouds” can be reduced.
\end{quoting}

\textbf{Response 9.} The Introduction now states once (p. 2): “(The scanner also records panoramic RGB imagery that colorizes the point cloud; the planning and evaluation in this paper concern the geometric measurements only.)” Repeated “colorized” phrases have been removed throughout; apart from that definitional sentence, the word now appears only in the caption of the colorized rendering in Figure 15, where color is genuinely shown.

\section*{Comment 10 (Minor) — Long sentences / English}

\begin{quoting}
The long sentence on page 1, lines 25–29 should be divided into shorter sentences. The English should be carefully revised throughout the manuscript, including the long captions and sentences.
\end{quoting}

\textbf{Response 10.} The page-1 sentence has been split (“…is growing. A dense and accurate point cloud is becoming an infrastructural asset: downstream robots localize against it, plan collision-free motion within it, and reason over it.”). We have additionally shortened the longest captions (Figures 2, 3, 7, 8, 11, 13–15; Tables 2–3) and made a careful English pass over the manuscript, which also addresses the “Quality of English Language” item of the review form.

\section*{Comment 11 (Minor) — Citation after author names}

\begin{quoting}
When reference authors are named in the text, the citation should follow their names directly. For example: “Mozaffar and Varshosaz [6]…”
\end{quoting}

\textbf{Response 11.} Fixed (Section 2, p. 5): “Mozaffar and Varshosaz [6] optimize scanner placement to reduce occlusions, and Prieto et al. [5] drive a next-best-scan policy…”. The same convention is applied to the newly added named citations (Dehbi et al., Knechtel et al.).

\section*{Comment 12 (Minor) — Abbreviations at first use}

\begin{quoting}
Define every abbreviation at first use, including TLS, SLAM, LOS, NBV, and MPPI, even if an abbreviation list is provided later.
\end{quoting}

\textbf{Response 12.} All abbreviations are now defined at first use: 3D, TLS, LiDAR, SLAM, GPR (Section 1); LOS (Section 1); NBV (Section 2.1); BIM (Section 2.3); MPPI (Section 5.1). The abbreviations list has been extended accordingly (LiDAR, GPR, BIM added; p. 25).

\section*{Comment 13 (Minor) — Figure citation order}

\begin{quoting}
Figure 4 is cited before Figures 1–3 are introduced. Figures should be numbered and cited in the order of first appearance.
\end{quoting}

\textbf{Response 13.} Fixed. The coverage-model figure (previously Figure 4, cited in Section 3 before Figures 1–3 appeared) has been moved to Section 3.2, where it is first cited; it is now Figure 1 (pp. 6–7), and all figures are numbered in order of first citation.

\section*{Comment 14 (Minor) — Figure 1 (architecture) label size}

\begin{quoting}
The labels in Figure 1 are too small at the normal manuscript scale. Please enlarge the text and simplify the diagram where possible. The caption should also be shortened, with detailed explanations moved into the main text.
\end{quoting}

\textbf{Response 14.} The architecture diagram (now Figure 2, p. 11) has been redrawn: labels are roughly twice as large at manuscript scale, the layout is a compact two-row pipeline, the redundant color legend was removed (the containers are labeled directly), and the caption has been shortened to two sentences with the component description kept in the Section 5.1 text.

\section*{Comment 15 (Minor) — Test room described in caption only}

\begin{quoting}
The industrial test room and its relevant characteristics are described mainly in the Figure 3 caption. The environment should first be introduced in the main text, followed by a citation to the figure.
\end{quoting}

\textbf{Response 15.} The test room is now introduced in the main text (Section 5.2, p. 12): open exhibition/staging space, reflective floor, display wall of monitors and product panels, freestanding equipment cases and furniture creating self-occlusion, floor markings for the survey reference layout, and its approximate extent of 16 × 8.5 m (about 105 m² of mapped free floor area). The figure caption has been reduced to a single sentence.

\section*{Comment 16 (Minor) — Millimeter units}

\begin{quoting}
Report millimeter-scale registration errors in millimeters rather than meters. For example: (0.004–0.005 m) should be written as 4–5 mm; (0.003–0.006 m) should be written as 3–6 mm. The units should be consistent throughout the manuscript.
\end{quoting}

\textbf{Response 16.} All registration errors are now reported in millimeters (Sections 6.7, pp. 21–22): bundle errors 4–5 mm, per-link errors 3–6 mm, disk single-link 2 mm, quality band ≤ 15 mm. Table 3’s column is now “Bundle error (mm)”. Units are consistent throughout the text, tables, and abstract.

\section*{Comment 17 (Minor) — Repetition}

\begin{quoting}
Several explanations are repeated in the abstract, introduction, results, discussion, figure captions, and conclusions, particularly those concerning: colorized point clouds; survey-grade registration; single-link and multi-link registration networks; and the difference between disk and visibility coverage. The manuscript would benefit from removing this repetition.
\end{quoting}

\textbf{Response 17.} We removed repeated explanations of colorized point clouds (now defined once), the single-link/multi-link registration contrast (explained once in Section 6.7 and only referenced elsewhere), and the disk-vs-visibility difference (defined in Sections 1 and 3 and thereafter referenced); shortened captions no longer duplicate the surrounding text.

\section*{Comment 19 (Minor) — Caption length}

\begin{quoting}
Many figure and table captions are unnecessarily long. Please shorten them and move interpretation or methodological detail into the main text.
\end{quoting}

\textbf{Response 19.} Captions of Figures 2, 3, 7, 8, 11, 13, 14, 15 and Tables 2, 3 have been shortened, with interpretation and methodological detail moved into the main text (Sections 5.1, 5.2, 6.2, 6.5–6.7).
